% Threat model and formal problem statement (task 10.2)
% Grounded in: README.md, related_work.md, data/_provenance/, the attack
% categories already present in eval/results/evaluation_raw.csv (injection,
% authority, corpus, perplexity), and the adversarial/credential-exposure
% findings in eval/results/summary_report.md.

\section{Threat Model and Problem Statement}
\label{sec:threat-model}

\subsection{System Model}
We consider a Retrieval-Augmented Generation (RAG) pipeline consisting of a
document corpus $C$, a retriever $R$, and a downstream language model $M$.
At query time, $R$ selects a set of documents from $C$ and passes them to
$M$ as context. PDS sits between $R$ and $M$ (equivalently, it can be run
offline against $C$ itself): each candidate document $d$ is scored before
it is allowed to influence $M$'s output.

\subsection{Attacker Capabilities}
We assume an attacker who can introduce or modify at least one document
that is later retrieved into $C$ -- for example, by contributing content to
a corpus with open or loosely-moderated ingestion (a wiki, a shared
documentation repository, a user-submitted knowledge base), or by having
one of their pages indexed by a web-scraped source. Consistent with prior
work on practical RAG poisoning attacks requiring only a single injected
document rather than large-scale corpus control~\cite{zhang2025practical},
we do \emph{not} assume the attacker needs to control a large fraction of
the corpus: our adversarial evaluation set (\texttt{data/adversarial/})
specifically tests single-document, hand-crafted attacks.

The attacker's goal is to have $M$ follow an injected instruction, leak
information, or otherwise deviate from its intended behavior as a result of
processing the planted document.

\textbf{The attacker CANNOT:}
\begin{itemize}
  \item Modify the retriever's ranking algorithm or the language model's
    weights, prompts, or system instructions directly.
  \item Observe PDS's internal detector scores or decision boundaries at
    attack-construction time (black-box with respect to the deployed
    detector, though we do evaluate white-box-style adversarial attempts
    in \texttt{data/adversarial/}, where the attacker is assumed to know
    the \emph{general approach} -- pattern/embedding/perplexity scoring --
    but not the exact calibration bounds or thresholds).
  \item Guarantee retrieval: PDS assumes the defender controls the
    detection step, not whether a malicious document gets retrieved in the
    first place (that is the retriever's / corpus-governance's
    responsibility, out of scope here).
\end{itemize}

\subsection{Attack Surface / Categories}
Our evaluation dataset covers four primary categories, matching the labels
used throughout \texttt{eval/results/}:
\begin{enumerate}
  \item \textbf{Instruction injection} -- explicit override phrases
    (\emph{"ignore previous instructions"}, fake system tags).
  \item \textbf{Authority spoofing} -- content that impersonates a
    legitimate authority (a fake security policy, a fake administrator
    notice) to lend false credibility to an instruction.
  \item \textbf{Corpus poisoning} -- semantically out-of-place content
    inserted into an otherwise trusted corpus, without necessarily using
    explicit override language.
  \item \textbf{Perplexity artifacts} -- text that is unnaturally smooth
    (often a sign of adversarially optimized or template-generated attack
    text) or unnaturally garbled.
\end{enumerate}
Our adversarial red-teaming set adds a fifth, informally-named category
surfaced during evaluation: \textbf{credential exposure framed as
boilerplate} -- a malicious instruction (e.g., shipping default
credentials) phrased as neutral, generic documentation/license text
specifically to dilute embedding-anomaly and avoid pattern triggers. See
Section~\ref{sec:limitations} for how this was handled.

\subsection{Defender Assumptions}
PDS assumes the defender:
\begin{itemize}
  \item Has access to full document text before the trust decision is made
    (i.e., detection happens pre-retrieval-trust, not by inspecting $M$'s
    output after the fact).
  \item Has a small corpus of known-clean reference documents
    (\texttt{data/clean/}) to calibrate the embedding-anomaly detector
    against.
  \item Operates under a latency budget: detection must not make the RAG
    pipeline unusably slow, which motivates the staged/tiered cascade
    design (Section~\ref{sec:limitations} notes the latency measurement
    caveats).
  \item Does \emph{not} have ground-truth poisoned/clean labels at
    deployment time -- only at calibration time, using the fixed,
    pre-computed bounds described in \texttt{eval/calibrated\_fusion.py}.
\end{itemize}

\subsection{Formal Problem Statement}
For a document $d$, each detector $i \in \{\text{pattern}, \text{embedding},
\text{perplexity}\}$ produces a score $s_i(d) \in [0, 1]$ and a latency cost
$\ell_i(d)$. The fusion layer maps the score vector to a verdict
$v(d) \in \{\text{Allow}, \text{Flag for review}, \text{Reject}\}$.

We frame the staged pipeline as a sequential decision / early-exit cascade:
at each stage $i$, the defender either (a) stops and emits a verdict based
on scores collected so far, or (b) pays cost $\ell_{i+1}(d)$ to run the next
detector. The objective is to choose a stopping rule $\pi$ minimizing
expected total cost:
\[
  \mathbb{E}_{d \sim C}\Big[\, \alpha \cdot \mathbb{1}[\text{false negative}]
  + \beta \cdot \mathbb{1}[\text{false positive}]
  + \gamma \cdot \textstyle\sum_{i \in \pi(d)} \ell_i(d) \,\Big]
\]
where $\alpha, \beta, \gamma$ weight missed-detection cost, false-alarm
cost, and latency cost respectively, and $\pi(d)$ is the set of stages
actually run for document $d$ under policy $\pi$. The fixed-order staged
pipeline in \texttt{eval/staged\_pipeline.py} is a simple, hand-chosen
instance of $\pi$ (run cheap, high-precision detectors first; only pay for
perplexity when inconclusive); Phase 5's cost-aware cascade
(\texttt{fusion/cost\_aware.py}) is intended to choose $\pi$'s stage order
and early-exit thresholds more systematically.

\subsection{Explicitly Out of Scope}
\begin{itemize}
  \item Attacks on the retriever's ranking function itself (SEO-style
    manipulation to win retrieval, independent of document content).
  \item Attacks against the language model $M$ after a document has
    already been allowed through (jailbreaks that don't rely on retrieved
    content).
  \item Colluding multi-document attacks, where no single document is
    individually suspicious (planned future work, Phase 8.3).
  \item Adaptive attackers with white-box access to PDS's exact calibration
    bounds and thresholds (our adversarial set assumes attacker knowledge
    of the \emph{general} detection approach, not the specific deployed
    parameters).
\end{itemize}
